% !TEX root = ../main.tex

\chapter{绪论}
随着区块链技术与智能合约的迅猛发展，去中心化应用已在金融、数字资产等领域得到广泛部署\cite{10.1145/3716323}。然而，智能合约的不可篡改性使得程序漏洞的代价日益高昂。在此背景下，专业安全审计公司应运而生，通过程序验证技术手段对合约代码进行安全性审查\cite{VIDAL2024112160}。然而，传统验证范式面临固有矛盾：审计过程要求完整暴露程序逻辑，使审计公司积累的核心验证技术与知识产权难以得到有效保护\cite{mota2023formallyverifyingrealworld}。
零知识证明(ZKP)技术能够在不泄露任何额外信息的前提下验证陈述的正确性\cite{sheybani2025zeroknowledgeproofframeworkssystematic}，天然契合隐私保护验证的需求。通过将程序验证与零知识证明技术相结合，可以构建一种新型验证范式：在证明程序满足安全规范的同时，严格保护程序实现细节与验证方法的机密性。\cite{10.1007/978-3-030-45721-1_24}这种方法不仅能够保障审计双方的知识产权安全，还可生成可复用的简洁证明，使第三方能够以极低的计算开销验证程序的正确性，从而有效消除重复验证带来的冗余计算，为隐私敏感场景下的程序验证提供切实可行的技术方案。

\section{研究背景}
程序验证（Program Verification）是通过形式化方法确保软件满足规范要求的关键技术，在金融、自动驾驶、智能合约等安全关键系统中具有重要应用价值\cite{10.1145/3605098.3636008, 10.1007/978-3-030-65277-7_5, BORCHERS2025103256}。随着区块链技术的发展，程序验证的重要性愈发凸显，智能合约的不可篡改性使得程序漏洞代价高昂，2022年至2024年间因智能合约漏洞导致的经济损失已达数十亿美元\cite{gogol2024sokdecentralizedfinancedefi}。目前已发展出多种程序验证技术，包括通过符号值探索所有可能执行路径\cite{10.1007/978-3-031-76554-4_13}的符号执行（Symbolic Execution）、在有限执行步数内验证性质\cite{10.1007/978-3-540-45069-6_2}的有界模型检测（Bounded Model Checking），以及自动发现循环不变式和函数前后置条件\cite{beyer2024augmentinginterpolationbasedmodelchecking}的不变式生成（Invariant Generation）。这些技术通常将程序验证问题归约为位向量理论上的可满足性模理论（SMT）问题，并通过位展开（Bit-blasting）技术进一步转换为底层SAT实例\cite{10.1145/3597926.3598034}。对于生成的SAT公式，底层SAT求解器若判定其可满足，则产生使公式为真的赋值；否则生成分辨证明（Resolution Proof）证明不存在这样的赋值。然而，该验证范式面临两个关键挑战：首先，证明验证正确性通常需要披露程序及其对应的SAT公式与证明，从而暴露具体实现细节并带来知识产权风险；其次，多方独立重复验证导致冗余计算开销，而盲目信任又削弱了验证的可信度。例如，金融科技公司若要向监管机构证明其交易算法满足特定性质，必须暴露SAT公式和程序本身，从而在透明性与机密性之间面临不可接受的权衡。

零知识证明（Zero-Knowledge Proof, ZKP）技术为解决上述困境提供了理论基础与可行路径。零知识证明是由Goldwasser、Micali和Rackoff于1985年提出的密码学协议\cite{DBLP:journals/siamcomp/GoldwasserMR89}，该协议允许证明者在不透露任何额外信息的情况下，向验证者证明某个陈述的真实性。这一特性天然契合隐私保护程序验证的核心需求。特别是零知识简洁非交互式知识论证（Zero-Knowledge Succinct Non-Interactive Argument of Knowledge, zk-SNARK）的发展，通过将计算问题编码为算术电路或多项式约束，实现了证明规模与验证时间的\cite{groth16}，为大规模程序验证提供了切实可行的技术路径。零知识证明系统能够生成紧凑的、可重用的证明，使得验证方能够以极低的计算开销验证计算的正确性，同时可证明地隐藏被验证程序的内部逻辑，从而有效解决隐私暴露与次线性增长冗余计算的双重挑战。\cite{10.5555/2671225.2671275}

近年来，零知识虚拟机（Zero-Knowledge Virtual Machine, zkVM）与普朗克（PLONKish）算术化方案的技术突破进一步拓展了零知识证明的应用边界\cite{DBLP:journals/iacr/GabizonWC19}。zkVM通过在虚拟机层面集成零知识证明能力，使得开发者能够以通用编程语言编写程序并自动生成相应的零知识证明，极大地降低了零知识应用的开发门槛。PLONKish算术化方案则通过引入自定义门、查找表等机制，显著提升了电路表达能力与证明效率。这些技术进展为构建高效的隐私保护程序验证系统奠定了坚实基础，使得在保护程序实现细节与验证方法机密性的同时，实现高效、可复用的程序正确性证明成为可能。\cite{Gabizon2019PLONKPO}

\section{国内外研究现状}
在程序验证与零知识证明交叉领域，当前的研究主要聚焦于如何利用零知识证明技术实现隐私保护的程序正确性验证。随着区块链技术的广泛应用与智能合约安全需求的日益增长，在不暴露验证细节的前提下证明程序满足特定安全属性\cite{10.1007/978-3-031-82703-7_11}，已成为学术界与工业界共同关注的重要研究方向。然而，现有验证方法通常要求完整暴露程序，这不仅使被验证方的商业机密面临泄露风险，也使验证方的核心技术难以得到有效保护。为应对这一挑战，研究者们提出了多种技术路径与解决方案。


\subsection{基于简洁非交互式证明的C程序执行验证}
Ben-Sasson等人在CRYPTO 2013上提出了SNARKs for C系统\cite{10.1007/978-3-642-40084-1_6}，实现了对C程序是否正确执行的简洁零知识证明。该系统将程序执行问题转化为算术电路可满足性问题，从而利用零知识证明技术进行高效验证。系统的核心技术贡献体现在三个方面：首先，设计了TinyRAM: 一种面向高效验证优化的极简非确定性随机存取机架构，其指令集经过精心设计，使得状态转移函数能够被编码为规模较小的算术电路；其次，开发了从GCC到TinyRAM的编译器后端，使得开发者能够直接使用C语言子集编写待验证程序；第三，构建了高效的零知识线性概率可检证明（Probabilistically Checkable Proof, PCP），并通过优化转换方法将其转换为对应的zk-SNARK。

在技术实现上，系统充分利用非确定性来优化电路规模。给定程序$P$和时间界$T$，系统生成的电路规模仅为$\tilde{O}(|P| \cdot T)$，相比现有电路生成器产生的二次规模电路有显著改进。证明仅包含12个群元素（共2576比特），且证明规模与电路大小和输入无关，验证仅需$O(|x|)$次密码学运算。

然而，该系统存在若干局限：首先，依赖于可信设置（Trusted Setup），需要可信第三方生成证明密钥和验证密钥；其次，证明生成的计算开销仍然较高；第三，系统主要关注程序执行的正确性证明，而非程序不满足某些性质的证明。

\subsection{基于零知识证明的静态程序分析}
针对静态分析过程中的程序隐私保护需求，Fang等人在CCS 2021上首次提出了零知识静态程序分析（Zero-Knowledge Static Program Analysis）的概念\cite{10.1145/3460120.3484795}。静态程序分析工具能够自动证明程序的诸多有用属性，然而，若需向第三方证明程序满足某一属性，传统方法必须暴露程序源代码。零知识静态分析的核心目标在于：证明者能够构造关于静态分析结论的零知识证明，在向验证者证明分析结果正确性的同时，完全隐藏程序本身。

该研究的关键技术贡献体现在以下方面：首先，研究团队针对过程内（Intra-procedural）和过程间（Inter-procedural）抽象解释提出了新颖的零知识证明方案。抽象解释是静态分析的核心技术框架，通过在抽象域上计算不动点来推断程序属性。然而，将抽象解释算法直接转换为零知识证明电路会导致极高的计算开销。为此，研究团队设计了专门针对抽象解释结构特点的优化方案，相比使用现有通用方案直接编码对应的静态分析算法，证明效率获得了显著提升。

在技术实现上，该系统支持两种零知识证明后端。第一种是基于配对（Pairing-based）的零知识证明方案，具有证明规模极小、验证速度极快的优势。实验结果表明，对于2,000行规模的程序，系统能够在1,738秒内完成控制流分析的零知识证明生成，所生成的证明仅为128字节，验证时间仅需1.4毫秒。第二种是基于离散对数（Discrete-log）的透明零知识证明方案，无需可信设置。在该方案下，系统能够对12,800行程序的污点分析在406秒内生成证明，证明规模为282千字节，验证时间为66秒。

实验评估在真实程序和合成程序上进行了系统性验证，结果表明该方法在实际程序规模上具有可行性。控制流分析（Control Flow Analysis）和污点分析（Taint Analysis）作为两种典型的静态分析技术，分别验证了系统在不同分析类型上的有效性。控制流分析能够追踪程序中函数调用的可能目标，而污点分析能够追踪敏感数据在程序中的传播路径，两者均为安全验证中的常用技术。

然而，零知识静态程序分析也存在若干局限：其一，该方法基于抽象解释框架，主要适用于可以通过抽象域上的不动点计算来表达的程序属性，对于需要精确推理的验证任务（如基于SAT/SMT求解的程序验证）支持有限；其二，虽然相比朴素方法有显著改进，但证明生成的计算开销仍然可观，限制了其在超大规模程序上的应用。

\subsection{基于交互式零知识证明的布尔公式不可满足性验证}
针对程序验证中的UNSAT问题，Luo等人在CCS 2022上提出了ZKUNSAT系统\cite{10.1145/3548606.3559373}，首次实现了在零知识条件下高效证明布尔公式的不可满足性。其核心洞察在于：虽然对于一般的UNSAT公式，归结证明（Resolution Proof）并不总是存在，但对于源自程序验证等实际应用的编码问题，现有SAT求解器通常能够高效地生成归结证明。

ZKUNSAT的关键技术贡献体现在三个方面：首先，该系统仅对归结证明的推导过程进行零知识建模，而非直接证明SAT公式本身。这一设计使得系统与现有SAT求解器完全解耦——在获得SAT公式后，证明者可以使用任意现有的SAT求解器（如PicoSAT、MiniSat等）独立求解并生成归结证明，无需修改底层求解器实现。这种模块化设计不仅保持了与现有验证工具链的完全兼容性，还使得系统能够受益于SAT求解器领域的持续进步。

其次，ZKUNSAT采用了创新性的代数编码方案，将公式子句表示为有限域上的低次多项式。具体而言，每个子句被编码为一个多项式，其根对应于使该子句为假的变量赋值。该编码方案充分利用了归结证明的结构特性：在实践中，归结证明通常具有较小的宽度（即每个子句包含的文字数量较少）。通过同时暴露证明的长度和宽度，系统实现了显著优化的验证协议。归结步骤的正确性验证被归约为多项式关系的检验，从而可以利用现有的零知识多项式承诺方案高效实现。

第三，该系统结合零知识随机访问论证（ZK Random-Access Argument）实现高效的证明验证。由于归结证明的验证需要频繁访问之前推导出的子句，研究团队采用了FlexZKArray功能来支持零知识条件下的随机内存访问，从而避免了朴素实现中$O(n^2)$的验证复杂度。

实验评估表明，ZKUNSAT在实际程序验证基准测试上展现出良好的性能。研究团队使用CBMC（C Bounded Model Checker）将C程序编码为布尔公式，并在SV-COMP基准测试集上进行了系统性评估，包括Linux内核驱动程序、Windows NT设备驱动程序和Intel密码学模块等实际验证任务。实验结果表明，ZKUNSAT能够在合理时间内验证中等长度、宽度高达2800的归结证明。

然而，ZKUNSAT存在两个主要局限：其一，作为交互式协议，证明者与验证者之间需要多轮通信，这在某些应用场景中可能带来额外的通信开销，且不适用于需要非交互式证明的区块链等场景；其二，该系统仅对Resoltuion Proof过程进行了零知识化，SAT公式对验证者公开，进而验证者可以通过SAT Solver对公式进行求解，未覆盖程序验证的全流程。


综上所述，当前零知识证明与程序验证的交叉研究已取得显著进展：SNARKs for C实现了C程序是否执行的非交互式零知识验证；零知识静态程序分析证明了程序具备某种属性而隐藏程序细节；ZKUNSAT解决了布尔公式不可满足性的高效零知识证明问题。然而，现有工作仍存在若干局限：SNARKs for C关注程序执行而非程序属性验证；ZKUNSAT采用交互式协议且仅支持纯布尔公式\cite{10.1145/3548606.3559373}；零知识程序静态分析证明生成的计算开销较大，限制了其真实场景的应用。更重要的是，现有方法尚未充分整合zkVM技术与高效的PLONKish证明系统来实现整个验证过程中实现对程序语义信息的端到端隐私保护与验证技术的保护。因此，如何构建一种非交互式、高效且同时支持SAT/UNSAT验证的隐私保护程序验证框架，仍是该领域面临的重要挑战。

\section{研究内容}

本文提出了一种基于零知识证明的两阶段协议，利用零知识虚拟机（zkVM）和PLONKish算术电路的最新技术进展构建零知识证明，解决隐私保护程序验证的挑战并减少冗余计算开销。第一阶段通过自定义HyperPlonk协议增强零知识虚拟机，执行两类程序验证任务：一是基于SAT的验证，将公共程序属性和私有程序实现通过静态单赋值变换转换为布尔可满足性（SAT）公式；二是基于Symbolic Execution的验证，通过探索程序路径生成路径约束（SMT公式），进而通过位分解技术转换为SAT公式，用于可达性验证和漏洞检测。两种验证方式的SAT编码过程均在零知识虚拟机中执行，避免程序代码段、控制流结构和中间状态的暴露。第二阶段利用基于PLONKish的算术电路验证SAT求解结果，针对可满足（SAT）和不可满足（UNSAT）两种情况分别构建PLONKish电路—对不可满足实例编码归结证明，对可满足实例验证满足赋值，两者都在整个验证过程中隐藏SAT公式和求解细节。阶段之间通过SAT公式的多项式承诺进行密码学链接，确保公式表示在阶段间的一致性。本文的主要贡献如下：
\begin{itemize}
    \item 提出了一种基于零知识虚拟机的隐私保护程序验证方法，通过引入自定义HyperPlonk协议进行增强，支持SSA到SAT编码和符号执行两种验证方案，在保持程序隐私的同时实现了相比原生zkVM三倍的性能提升。
    \item 设计了一种新颖的PLONKish算术电路，针对SAT和UNSAT两种情况分别采用不同策略：对UNSAT实例编码归结证明，实现验证时间与子句宽度无关；对SAT实例进行满足赋值验证，实现验证时间与公式规模无关。两种方法均能生成高效可验证的隐私保护证明。
    \item 在真实程序验证基准测试集上进行了系统性实验评估，验证了本文方法的实用性和有效性。
\end{itemize}



\section{本文章节结构}
本文的组织结构如下：

第二章介绍了本文研究所需的理论基础和技术背景。主要内容包括零知识证明的基本概念与性质，非交互式零知识证明（NIZK）及zk-SNARK协议，Plonk算术化方案的核心原理，零知识虚拟机（zkVM）的设计与实现，程序验证中的符号执行技术，以及布尔可满足性问题（SAT）的归约方法。

第三章详细阐述阶段一中基于零知识虚拟机的程序验证实现。首先介绍zkWASM虚拟机的基础架构，然后详细描述SSA程序到SAT公式的转换算法，包括布尔逻辑编码规则、嵌套条件表达式处理和Tseitin变换；接着介绍基于符号执行的程序分析方法，包括路径探索、约束生成和SMT到SAT的位分解转换；最后说明HyperPlonk协议对zkVM的性能增强机制。

第四章深入阐述阶段二中基于Plonkish算术电路的SAT求解验证系统。首先介绍halo2电路系统的基本架构与约束类型；然后针对UNSAT实例，详细描述归结证明的挑战值编码方案、矩阵构造算法和约束系统建立，证明验证时间与子句宽度无关的特性；接着针对SAT实例，设计双矩阵赋值存储方案和互斥性约束，证明验证时间与公式规模无关的特性；最后通过定理证明两种方案的完备性、可靠性和零知识性。

第五章展示了本文协议的实现与实验评估。详细介绍实验环境配置和基准测试集的选取，通过在Boolector等标准测试集上的系统性评估，验证了阶段一相比原生zkVM实现三倍性能提升，以及阶段二UNSAT和SAT验证的高效性。实验结果表明本文方法在保持零知识性的同时实现了高效的程序验证。

第六章总结了本文的主要工作和研究成果，分析了研究方法的创新点与局限性，并对未来的研究方向提出了展望。

文末致谢各位师友的指导与帮助，并列出了本研究所参考的相关文献。